\section*{Chapter \ref{ch:iv:cpp} --- C++}

\begin{solution}{iq:iv:cpp:1}
\texttt{Base}, then \texttt{Derived}, then \texttt{Derived}. While \texttt{Base}'s constructor runs, the \texttt{Derived} part does not
exist yet and the object's dynamic type is \texttt{Base}, so the virtual call resolves to \texttt{Base::name}; in
\texttt{Derived}'s constructor and afterwards through the reference it resolves to \texttt{Derived::name}. The same holds in
destructors, in reverse. The chapter's test compiles and runs the snippet.
\iqlookfor{the dynamic type during construction, and the rule's mirror in destruction.}
\end{solution}

\begin{solution}{iq:iv:cpp:2}
\texttt{make second}, \texttt{make first} (members in declaration order), \texttt{body}, \texttt{make temp}, \texttt{drop temp}
(the temporary dies at the end of the full expression, after \texttt{use} returns), \texttt{after t}, then at the end of
\texttt{main} \texttt{drop first}, \texttt{drop second} (reverse order of construction).
\iqlookfor{declaration order, temporary lifetime to the end of the full expression, reverse destruction.}
\end{solution}

\begin{solution}{iq:iv:cpp:3}
All five: destructor (\texttt{delete[]}), copy constructor and copy assignment (deep copy, or delete them), move constructor and move
assignment (steal the pointer, leave the source empty, \texttt{noexcept}). Better, the rule of zero: hold a
\texttt{std::vector<T>} or a \texttt{std::unique\_ptr<T[]>} and declare none of them; the compiler generates correct ones.
\iqlookfor{the rule of five and the preference for the rule of zero.}
\end{solution}

\begin{solution}{iq:iv:cpp:4}
Resource acquisition is initialisation: a resource is owned by an object whose constructor acquires it and whose destructor
releases it, so release happens on every exit path, exceptions included. Examples: a \texttt{std::lock\_guard} around a shared
book structure; a session object that logs out and closes the socket in its destructor; a memory-mapped file or a pinned memory
region released on scope exit.
\iqlookfor{the definition in terms of lifetime and exception safety, with concrete examples.}
\end{solution}

\begin{solution}{iq:iv:cpp:5}
\texttt{unique\_ptr} with the default deleter is one pointer, and moving it costs a pointer copy. \texttt{shared\_ptr} is two
pointers (object and control block) plus a separately allocated control block (unless made with \texttt{make\_shared}), and every
copy and destruction updates an atomic reference count, which is contended across threads. The snippet prints \texttt{8 16 32}:
those are libstdc++ sizes on x86-64 (implementation-defined); the standard guarantees only the interfaces. Use
\texttt{unique\_ptr} unless ownership is genuinely shared.
\iqlookfor{the costs, and the distinction between guaranteed and implementation-defined facts.}
\end{solution}

\begin{solution}{iq:iv:cpp:6}
\texttt{lvalue}, \texttt{rvalue}, \texttt{rvalue}, \texttt{lvalue}. Inside \texttt{g}, \texttt{s} has a name and so is an lvalue even
though its type is an rvalue reference. Write \texttt{f(std::move(s))} (or \texttt{std::forward} in a template taking a forwarding
reference).
\iqlookfor{value category of a named rvalue reference, and \texttt{std::move}.}
\end{solution}

\begin{solution}{iq:iv:cpp:7}
\texttt{symbol()} returns a temporary \texttt{std::string} that is destroyed at the end of the declaration, so the
\texttt{string\_view} dangles and reading \texttt{s[0]} is undefined behaviour. (A short string might sit in the string's internal
buffer and appear to work, which hides the bug; the snippet uses a long one.) AddressSanitizer reports a heap-use-after-free,
as the chapter's test asserts. Fix: keep a \texttt{std::string} (\texttt{auto s = symbol();}). To find such bugs: build tests with
sanitizers, and enable the compiler's dangling warnings where available.
\iqlookfor{temporary lifetime and views, and sanitizers as the detection method.}
\end{solution}

\begin{solution}{iq:iv:cpp:8}
Nothing can be predicted: \texttt{INT\_MAX * 2} overflows a signed \texttt{int}, which is undefined behaviour; the chapter's test
asserts only that UndefinedBehaviorSanitizer reports ``signed integer overflow''. Because overflow cannot happen in a valid
program, the compiler may assume \texttt{i + 1 > i} is always true and turn a loop bounded by it into an infinite loop, or remove
an overflow check written as \texttt{if (a + b < a)}. Use unsigned arithmetic when wrap-around is wanted, a wider type, or checked
arithmetic (\texttt{\_\_builtin\_add\_overflow}).
\iqlookfor{UB stated, not predicted, and what the optimiser does with it.}
\end{solution}

\begin{solution}{iq:iv:cpp:9}
\texttt{first} refers into the vector's buffer; the \texttt{push\_back}s reallocate it and the reference dangles, so reading it is
undefined behaviour (AddressSanitizer: heap-use-after-free). Fixes: \texttt{reserve(101)} before taking the reference, keep an index
instead of a reference, or copy the value.
\iqlookfor{reallocation invalidating references, and index-based or reserved alternatives.}
\end{solution}

\begin{solution}{iq:iv:cpp:10}
It prints \texttt{over limit} and \texttt{4294967295}. In \texttt{balance < limit} the \texttt{int} is converted to
\texttt{unsigned}, so $-1$ becomes $2^{32} - 1$ on this platform (32-bit \texttt{unsigned}, implementation-defined width) and the
comparison is false. A negative balance is reported as over a limit of 1, and the opposite mistake would let a breach through.
The compiler warns (\texttt{-Wsign-compare}); build with \texttt{-Werror}, and keep money and quantities in one signed type.
\iqlookfor{the usual arithmetic conversions and their consequence for a control.}
\end{solution}

\begin{solution}{iq:iv:cpp:11}
\omcode{code/interviews/22-cpp/cpp/iv_cpp.hpp}{13}{46}{A buffer with the rule of five; the moves leave the source empty and cannot
throw.}\label{lst:iv:cpp:buffer}
Moves must be \texttt{noexcept} so that \texttt{std::vector<Buffer>} moves elements when it grows (it copies if a move could throw,
to keep its strong guarantee). A moved-from \texttt{std::vector} is valid but unspecified; libstdc++ leaves it empty (the chapter's
snippet prints \texttt{v=0}), but portable code only destroys or reassigns it. The chapter's C++ test checks moves, copies and
self-assignment.
\iqlookfor{\texttt{std::exchange}-style moves, \texttt{noexcept} and why, and the moved-from contract.}
\end{solution}

\begin{solution}{iq:iv:cpp:12}
\omcode{code/interviews/22-cpp/cpp/iv_cpp.hpp}{48}{67}{A fixed-capacity ring buffer with monotonically increasing unsigned
counters.}\label{lst:iv:cpp:ring}
When full, \texttt{push} returns false and the caller chooses a policy: drop, block, or signal back-pressure (One Quant Book~13,
chapter~12). A power-of-two capacity replaces the modulo by a mask, and free-running unsigned counters make ``full'' and ``empty''
distinguishable without wasting a slot. The test compares it with a \texttt{std::deque} model over 100\,000 random operations.
\iqlookfor{the full/empty logic, the mask, and a stated full policy.}
\end{solution}

\begin{solution}{iq:iv:cpp:13}
\omcode{code/interviews/22-cpp/cpp/iv_cpp.hpp}{69}{88}{A constexpr table and a concept.}\label{lst:iv:cpp:constexpr}
If the loop multiplies once more, to $10^{19}$, the computation overflows \texttt{int64\_t}; in a constant expression undefined
behaviour is not allowed, so the compiler refuses to compile it (g++ reports ``overflow in constant expression''). That happened
while writing this chapter, and it is the argument for computing tables at compile time. The concept makes \texttt{notional} accept
any type with integer \texttt{price} and \texttt{qty} and reject others with a clear error; the test checks both with
\texttt{static\_assert}.
\iqlookfor{constexpr evaluation catching UB, and concepts as checked interfaces.}
\end{solution}

\begin{solution}{iq:iv:cpp:14}
\texttt{init b}, \texttt{init a}, \texttt{a=3 b=2}: within one translation unit, dynamic initialisation follows the order of
definition. Across translation units the order is unspecified, so \texttt{a} could be initialised before \texttt{b} and read
zero (static storage is zero-initialised first): the static initialisation order fiasco. Fix: make the dependency a function
returning a local static (initialised on first use, thread-safe), or make the values \texttt{constexpr}/\texttt{constinit} so that
they are initialised at compile time.
\iqlookfor{the within-unit guarantee, the cross-unit fiasco and the standard fixes.}
\end{solution}
